Deixe \(n\) ser um número inteiro positivo.
Um \emph{\(n\)-simplex} em \(\mathbb{R}^n\) é dado por \(n+1\) pontos \(P_0, P_1,\dots  , P_n\), chamados de seus vértices, que não pertencem todos ao mesmo hiperpano.
Para todo \(n\)-simplex \(\mathcal{S}\), denotamos por \(v(\mathcal{S})\) o volume de \(\mathcal{S}\), e escrevemos \(C(\mathcal{S})\) para o centro da esfera única que contém todos os vértices de \(\mathcal{S}\).

Suponha que \(P\) seja um ponto dentro de um \(n\)-simplex \(\mathcal{S}\).
Deixe \(\mathcal{S}_i\) ser o \(n\)-simplex obtido de \(\mathcal{S}\) substituindo seu \(i^{\text{th}}\) vértice por \(P\).
Prove que
\[ \sum_{j=0}^{n}v(\mathcal{S}_j)C(\mathcal{S}_j)=v(\mathcal{S})C(\mathcal{S}). \]
